\chapter[INTRODUCTION]{Introduction}\label{ch:I}

\section{Background of the Study}\label{sec:1.1}

The world is a huge place, and we may often find ourselves getting lost in
unfamiliar places. Throughout history, we have invented various tools, starting
with the compass and the map, to help us with navigation. Nowadays, navigation
is made easy thanks to modern tracking and mapping technology. In the
Philippines, applications like Waze help to ease the transportation
crisis~\cite{ref01}.

One of the most commonly used technologies is the Global Positioning System
(GPS), a publicly available mapping system created by the United States
government that uses satellites to approximate the location of a GPS receiver
contained in a device such as a laptop or a mobile phone. This allows users to
obtain directions to and from a particular place, and it is especially useful in
a travel setting. Popular applications that use GPS for navigation include
Google Maps, Waze, and Apple Maps.

However, due to its top-down view from the satellites, GPS technology works best
outdoors. It does not work well for buildings with multiple levels. To solve
this, we can use an indoor positioning system (IPS), which works like a
miniature version of the GPS. Specifically, we can install specialized tracking
devices in various places inside the building. Several technologies can be used
as positioning methods for the IPS, some of which include Wi-Fi, infrared,
Bluetooth, and computer vision.

To further improve the IPS in terms of visualization, we can implement an indoor
map where we can mark the locations of each landmark, as well as the location of
the user. These maps can be derived from the floor plans of the building, and
they can be implemented without an IPS. Examples of applications that allow us
to design indoor maps are Mappedin and ArcGIS Indoors. Google Maps also supports
indoor maps, but only a few buildings currently have this functionality.

\section{Statement of the Problem}\label{sec:1.2}

Globally, hundreds of thousands of buildings are being constructed every year.
In the Philippines, the total number of buildings constructed in the year 2022
was 169,432, which was 4.2 percent higher than the 162,619 constructions in
2021~\cite{ref02}. Since GPS is not an effective solution for mapping the inside
of these buildings, there is a growing need for indoor mapping solutions.

The Physical Sciences Building on the University of the Philippines Los Baños
campus is one of these buildings that need indoor mapping. New students and
visitors may find it difficult to navigate inside this building due to its
multiple floors and classrooms. Currently, there is a floor plan located in this
building, but it is not easy to find.

An IPS is a possible solution for mapping the Physical Sciences Building, but it
can be costly and challenging to build and maintain. It is also not very
scalable if the system were expanded to the entire campus, because it requires
external hardware for each floor in each building. Hence, we turn to indoor
mapping applications that can be accessed anywhere.

\section{Objectives of the Study}\label{sec:1.3}

The main objective of this study is to develop a web application that provides
an indoor map of the UPLB Physical Sciences Building. Specifically, it aims to:

\begin{enumerate}
\def\labelenumi{\arabic{enumi}.}
\item
  integrate point-to-point directions to aid users in finding the exact
  locations of classrooms, offices, and facilities within the UPLB Physical
  Sciences Building;
\item
  enable building administrators to generate a building's maps based on floor
  plans, allowing for updates and customization of the mapping content; and
\item
  assess the usability of the system through user testing to evaluate navigation
  accuracy, interface intuitiveness, and overall user satisfaction.
\end{enumerate}

\section{Significance of the Study}\label{sec:1.4}

The web application developed in this study will be beneficial to visitors and
new students who have classes inside the Physical Sciences Building. It will be
easier for them to find the classrooms that they need to get to, and it will
also save them time and prevent traffic in the corridors.

The results of this study may also serve as a basis for future indoor mapping
implementations in UPLB. Aside from determining the feasibility of such a
project, future studies will know what can be improved based on this study.

\section{Scope and Limitations of the Study}\label{sec:1.5}

The main focus of this research is to build a web application to navigate
through the inside of the Physical Sciences Building located at the University
of the Philippines Los Baños, Laguna. The end users of this research include
individuals who are unfamiliar with the building, regular occupants of the
building such as employees and students, and building administrators of the
Physical Sciences Building. Using the application requires access to network
connectivity, as the application may need to update map data and navigation
information as it is being used, and it is also required when logging in.

The application does not track the user's real-time location; instead, users
select their origin and destination, and the route is displayed on the map.
Accessible routes and emergency exit routes are not yet mapped. The evaluation
of the application was limited to six student users and four building
administrators from the University Planning and Maintenance Office (UPMO).

\section{Definition of Terms}\label{sec:1.6}

The following terms are used in this study as defined below.

\textbf{Administrator.} A user who logs in to the application and can create,
edit, and publish floor layouts through the layout builder, in addition to using
the navigation features.

\textbf{Draft layout.} An unpublished version of a floor layout that
administrators can save, load, rename, and delete before publishing it.

\textbf{Floor layout.} The digital representation of a floor of the building,
composed of rooms, hallways, stairs, bridges, and entrances positioned and sized
on a two-dimensional plane.

\textbf{Indoor positioning system.} A system that approximates a user's location
inside a building using devices installed in the building, such as Wi-Fi,
infrared, Bluetooth, or computer vision.

\textbf{Layout builder.} The administrator feature used to create, edit, and
publish the floor layouts on which navigation is based.

\textbf{Location-based services.} Technologies that use a device's location to
provide services such as object positioning, navigation, and mapping.

\textbf{Navigation graph.} The network of nodes (rooms, doors, hallway segments,
stairways, and the main entrance) and connecting edges that the application
builds from the published layout and on which it computes routes.

\textbf{Published layout.} The finalized layout that serves as the reference for
the map and routes shown to all users; also called the main layout.

\textbf{Regular user.} A student, employee, or visitor who uses the application
to search for locations and view routes within the building.

\textbf{Sandbox mode.} A separate workspace in which administrators can explore
the layout builder without affecting the main layout or the published map.

\textbf{System Usability Scale.} A ten-item questionnaire that produces a single
usability score from 0 to 100.

\textbf{Wayfinding.} The process of determining and following a route from an
origin to a destination.

\textbf{Wing.} A section of a floor of the Physical Sciences Building (Wing A,
Wing B, or Wing C) that users can select when viewing the map.
